% ============================================================
% MODULO 07 - LIMITACOES E EXPLICACOES ALTERNATIVAS
% ============================================================
\section{Limita\c{c}\~oes e explica\c{c}\~oes alternativas}

\subsection{Coer\^encia, n\~ao prova}

O percurso de M. repete, em escala individual, o que a interven\c{c}\~ao brasileira descreveu em grupo: a intensidade dos sintomas diminuiu, mas o tipo de TDAH permaneceu \cite{damasceno2025}. O relato tamb\'em se alinha ao achado de que espa\c{c}os verdes e abertos, somados ao conv\'ivio social, configuram o cen\'ario mais favor\'avel \cite{kuo2004}, e \`a revis\~ao que sustenta a plausibilidade do mecanismo restaurador \cite{hood2024}. H\'a ainda um contraponto digno de nota: a rotina de M. inverteu o que a literatura chama de ``d\'eficit de natureza'', fen\^omeno pelo qual a redu\c{c}\~ao do tempo livre ao ar livre na inf\^ancia contempor\^anea se associa a piores desfechos de desenvolvimento \cite{louv2015}. Na leitura humanista, a mudan\c{c}a de endere\c{c}o n\~ao ``tratou'' a crian\c{c}a: ampliou seu campo relacional. Aceita na rua, apesar da agita\c{c}\~ao, e aceita na sess\~ao, sem condi\c{c}\~oes, ela p\^ode atualizar uma pot\^encia que j\'a existia \cite{rogers2020}. O corpo que precisava interromper para existir passou a correr, negociar e descansar.

\subsection{Explica\c{c}\~oes alternativas}

Conv\'em considerar, sem ordem de import\^ancia: (a) matura\c{c}\~ao pr\'opria da faixa de 5,0 a 5,5 anos; (b) mudan\c{c}as de sono, rotina, escola e alimenta\c{c}\~ao; (c) al\'ivio do estresse familiar, capaz de influenciar o relato materno; (d) aprofundamento da alian\c{c}a terap\^eutica com o tempo; (e) vi\'es de expectativa -- a m\~ae deseja ver a filha bem na casa nova; (f) sazonalidade, com mais tempo ao ar livre no per\'iodo; (g) redu\c{c}\~ao de telas n\~ao mensurada; e, sobretudo, (h) medica\c{c}\~ao n\~ao declarada. Se algum psicof\'armaco foi iniciado, ajustado ou interrompido nesse intervalo, a leitura ambiental perde sustenta\c{c}\~ao. Sem SNAP-IV de linha de base, sem medida da dose de brincar e de telas e sem informa\c{c}\~ao farmacol\'ogica, este trabalho n\~ao afirma efic\'acia: afirma coer\^encia te\'orica e prop\~oe uma agenda.

\subsection{Limites do desenho}

Um caso \'unico, sem grupo de compara\c{c}\~ao e sem medidas padronizadas, n\~ao autoriza generaliza\c{c}\~oes. A converg\^encia entre o que a m\~ae conta e o que a terapeuta observa reduz -- n\~ao elimina -- o vi\'es de informante. A aus\^encia de dado sobre medica\c{c}\~ao permanece como a limita\c{c}\~ao mais grave e deve ser explicitada em qualquer vers\~ao a submeter. A verifica\c{c}\~ao emp\'irica da hip\'otese depender\'a de um desenho futuro com medidas repetidas: linha de base, seguimento em tr\^es meses e compara\c{c}\~ao entre contextos, conforme os crit\'erios de invalida\c{c}\~ao j\'a registrados na Metodologia.
